\chapter{37}



Der Abend hätte angespannt verlaufen können. Es kam in mehrfacher Hinsicht anders.
Die anfängliche Leichtigkeit stellte sich zusammen mit einem vorzüglichen Nachtessen
ein. Und sie war nicht zuletzt der Zurückhaltung der beiden zu
verdanken. Sobald Louis herumdruckste, liessen sie ihn. Ihre
Gelassenheit war wohltuend. Und er wusste, der Moment war günstig. Jetzt
hatte er die Gelegenheit, seine Brille loszuwerden. Mit verbogenem
Gestell und rausgesprungenen Gläsern lag die in seiner Hosentasche. Er
hatte keine andere Möglichkeit gesehen, als sie gewaltsam zu zertreten.
Wie oft hatte ihn die verzerrte Sicht geärgert, wäre er um ein Haar
hingefallen und hatten seine Augen getränt.
Joh brachte den Nachtisch. Und Louis legte den Rest Brille auf den
Tisch. Dazu machte er ein genervtes Gesicht. Manuela reagierte sofort.

«Hey, was ist denn mit deiner Brille passiert?»

Joh blickte ihn entsetzt an.

«Heute Morgen trat ich versehentlich drauf. Die lag neben meinem Bett,
ich stand auf und schwups, war sie flach.»

Johs Lösungsstrategien schienen ihn wohl auszuzeichnen. Doch dieser
Vorschlag kam ungelegen.

«In Brig kenne ich einen jungen Optiker bei Fielmann, Manuel, ein guter
Kerl,» Joh wandte sich an Manuela, «weisst du, Patricks Sohn, der ist neu
Filialleiter dort.»

Manuela brauchte einen Moment. Schliesslich nickte sie.
Louis musste viel Überzeugungskraft aufbieten, um den Vorschlag
abzuwenden. Sprach von einer ganz leichten Korrektur, und dass er
vorläufig ohne Brille zurecht käme. Und er gerne auf Johs Idee zurückkommen würde, sollte es denn wichtig werden. Das Thema schien erledigt. Und er hatte wieder freie Sicht. 

 In Louis hatten sie einen aufmerksamen Zuhörer. Er musste sich keine
Mühe geben. Ihre Geschichten erwiesen sich als optimale Ablenkung. Manuela sah zerknittert aus, fiel Louis auf. Ihre geröteten Wangen und
Augen zeugten von einem anstrengenden Tag. Langsam, pausenreich und
zwischen einzelnen Bissen beschrieb sie die Situation auf der Alp. Dabei
wandte sie sich immer wieder direkt Joh zu.

«Alles in allem,» schloss sie, «Mazi ist ein Held. Die Woche lief wie am
Schnürchen. Dafür, dass er alleine ist, einfach toll. Und zum Glück
gab’s keine Verletzungen. Auch kein Wolf in Sicht, hat er gesagt, na ja,
wie lange wohl?»

Manuela griff nach ihrem Glas und trank es leer. Joh hatte keine Zwischenfragen
gestellt. Er schien nachzudenken. Schliesslich schaute er zwischen
Manuela und Louis hin und her.

«Ich schlage vor, wir steigen erst übermorgen mit dir auf, Ciro. Mazi
trau ich einen letzten Tag alleine zu. So gut, wie’s oben läuft. Und du
kannst dich ausruhen, nach deinem anstrengenden Arbeitstag, Manuela.»

Manuela schien Johs Vorschlag zu gefallen. 

«Ihr hättet auch nur zu zweit hoch gehen müssen. Ich bin von
Davids Betreuer morgen zu einer Lagebesprechung in Brig eingeladen.
David ist auch dabei,» Joh atmete tief ein. Darauf neigte er sich zu
Louis.

«Für dich ist es bestimmt auch gut, Ciro, wenn du dich noch einen Tag in
Ruhe vorbereiten kannst. Ich geb dir entsprechende Lektüre und Manuela
ist da, wenn du Fragen hast. Was meint ihr?»

Johs Planänderung gab dem Abend augenblicklich eine neue Note. Sie
hatten Zeit. Kein frühes Aufstehen am nächsten Morgen. Manuela reckte
lächelnd den Hals und schaute zur Sofaecke hinter Joh. Daran angelehnt
hatte Louis seine Gitarre abgestellt.

«Du hast daran gedacht, Ciro.»

Sie sah überrascht und gleich wacher aus und zeigte auf die Gitarre.

«Ich bin sehr gespannt, was du für uns spielst.»

Louis würde tun können, was er gut konnte.

Der Raum, der an die offene Küche grenzte, war aussergewöhnlich weit.
Eigentlich machten Küche und Wohn-Esszimmer die gesamte Wohnfläche des
Erdgeschosses aus. Das stattliche Bauernhaus musste rigoros umgebaut
worden sein.
Manuela und Joh versanken in den tiefen Kissen des Sofas. Beide hatten
eine Bierflasche in der Hand. Louis griff sich seine \emph{Gibson}.
Während er auf dem unteren Teil der Sitzlandschaft nach einer passenden
Sitzposition suchte, begann er Passengers «Let her go» zu
beschreiben. Manuela meinte, den Song zu kennen und zu mögen. Louis
konnte sich nicht mehr zurückhalten. Es musste an der Stimmung,
diesem Frieden liegen, den Manuela und Joh ausstrahlten. Und auch daran,
dass Louis spürte, dass ihm ein bisschen Wahrheit ganz guttun würde. Er sprach von sich.

«Passengers Text spricht mich sehr an.»

Louis redete leise. Mit seinen Augen fixierte er seine linke Hand und
spielte einige Akkorde an. Seine Worte klangen wehmütig.

«Meine Freundin damals würdevoll zu verlassen, wäre vielleicht besser
gewesen. Aber ich konnte es einfach nicht. Noch heute 
komme ich kaum drüber hinweg. Eigentlich bin ich nach wie vor
schockiert.»

Joh zog die Augenbrauen hoch und nickte. Manuela rückte sich zurecht. Mit der einen Hand fasste Joh nach ihrer, mit der andern führte er die Bierflasche an die Lippen und nahm einen kräftigen Schluck.
Louis räusperte sich und begann schliesslich zu spielen. Nach den ersten Worten schaute er flüchtig hoch. Er bemerkte Manuelas glasige Augen. Sie wirkte auffallend angetan.

\qrblock{Music Travel Love (Passenger) \emph{«Let her go»}}{media/QR-Codes/031_Let_her_go_Music_Travel_Love_Spotify.png}

Joh klatschte, nachdem Louis die Gitarre abgestellt hatte. Dann stand er
auf und ging auf den Kühlschrank zu.

«Ciro, du bist unglaublich, wow,» sagte Manuela ergriffen und nahm ihr
Gesicht zwischen die Hände, «wie du Gitarre spielst und erst deine
Stimme, das war schön, danke.»

Joh hatte sich unterdessen neben Louis gestellt, blickte anerkennend zu
ihm hinunter und reichte ihm die Flasche.

«Ich glaube, da haben wir was verpasst. Du bist wohl ein bekannter
Musiker?»

Joh runzelte fragend die Stirn.

«Hm, nein, nein, noch nicht,» Louis lachte verhalten und suchte nach
Worten, viel mehr nach Inhalt, «ich spiele vorwiegend Cover-Versionen,
da kommst du nicht weit. Meine Sammlung an eigenen Songs ist noch am
Werden. Aber ich arbeite dran.»

Manuela stand auf.

«Ich finde deine Stimme wirklich einzigartig gut,» sie schaute ihn gross
an, «sehr schön.»

Sie blieb neben Joh stehen. Der legte den Arm um ihre Schultern, schob
seine Brille hoch und schien nachzudenken. Auch er wirkte sichtlich
beeindruckt.

«Live habe ich so was nicht oft gehört,» dann begann er laut zu lachen,
«schon gar nicht in unserem Wohnzimmer.»

Manuela löste sich sanft von Joh, gab ihm einen Klaps auf den Rücken und entfernte
sich Richtung Küche. Joh setzte sich.

«Spass beiseite, ich bin geflasht, Ciro. Doch wie du sagst, ich kann mir
gut vorstellen, dass du erst mit eigenen Produktionen etwas Grosses
erreichen kannst. Und wenn ich als Laie einen Blick auf diese
riesengrosse Musiklandschaft werfe, dann hat jemand, der es schafft,
wohl Erstaunliches kreiert.»

«Ja, das Musikbusiness ist hart, und ich bin überzeugt, dass auch Glück
dazu gehört.»

Manuela kam langsam und beidhändig beladen mit Desserttellern zum Sofa
zurück.

«Also, Ciro, mit \emph{einem} Song kommst du heute definitiv nicht
davon,» lachte sie, «aber erstmal gibt’s Nachtisch. Der Apfelkuchen
kommt direkt aus dem Ofen, ist noch warm.»

Manuela reichte Louis den Spritzsack mit geschlagener Sahne. Sie setzten
sich bequem in die Kissen. Wie bisher alles aus dieser Küche, schmeckte
der Kuchen grossartig. Die Ruhe während des Essens war angenehm.
Offensichtlich hatte er die beiden mit seiner Erklärung
zufriedengestellt.
Je wohler sich Louis fühlte, desto mehr interessierte er sich für
\emph{sie.} Für ihre Geschichten. Und nicht zuletzt für sie beide als
Paar. Die waren zehn Jahre zusammen und gaben sich wie junge
Frischverliebte? Wie mochte das gehen? Was war ihr Geheimnis?
Er durchbrach die Stille.

«Ihr habt eine gute Art miteinander umzugehen, find ich,» er schluckte 
ein weiteres Stück Kuchen hinunter, «ich kenne kaum Paare in eurem
Alter, die sich so,» Louis überlegte, «so
aufmerksam behandeln.»

Die beiden lachten. Er lauter als sie.

«Alte Verliebte, nenn es nur beim Namen, Ciro. Weisst du, als junger Mann
sah ich 60-jährige auch als alte Knacker.»

Johs Bemerkung schien Manuela zu belustigen. Sie kniff ihn unerwartet in
die Seite. Beinahe wäre ihm die Gabel aus der Hand gefallen. 

«Dennoch, so einfach ist es nicht,» Manuela sah wieder ernst aus, «wir
haben gemeinsam beschlossen, etwas dafür zu tun, um zusammen lebendig zu
bleiben.»

Sie hielt inne und blickte Joh an. Seine Augen verengten sich. Jetzt 
glaubte Louis wieder diesen ernsten, nachdenklichen Blick in Johs
Gesicht zu erkennen. Er genoss es, Joh beim Denken zuzusehen. Der Mann
erschien Louis stark. Und doch war er kein Ego-Kracher, wie er ältere
Männer oft erlebte. In Johs Nähe stellte sich für Louis ein Gefühl von
Sicherheit ein. Ganz ähnlich wie bei Grandpère, staunte er und freute
sich über den Vergleich.

«Manuela sagt’s, es ist nicht einfach, oder besser, es war nicht
einfach, jetzt finde ich,» und mit einem Blick zu ihr, «haben wir einen
Ausdruck von Zusammengehörigkeit und Liebe erreicht, der, der schwer in
Worte zu fassen ist.»

Joh schien mit sich zu ringen. Seine Worte hörten sich an wie eine
Liebeserklärung. Manuela nickte zustimmend. Als Zeuge, sozusagen mittendrin, spürte Louis, wie ihn
ein leichtes Gefühl von Peinlichkeit streifte. 
Louis schaute die beiden abwechselnd an.

«Und ja, Beziehungen brauchen Pflege und Hingabe. Und wie ich lernen
musste, auch den starken Willen, sich selbst immer wieder aufs Neue
kennenzulernen, zu wachsen, im besten Fall,» ergänzte Joh.

Louis erinnerte sich erneut an die ersten Stunden auf der \emph{Sunnewaid}. Wie
aussergewöhnlich Joh Manuela umsorgt hatte. Und daran, dass sich beide
für das Nachtessen umgezogen hatten. Dass irgend etwas Feierliches in
der Luft gelegen hatte, das Louis schwergefallen war, einzuordnen. Er
hüstelte.

«Hatte der gestrige Abend eine besondere Bedeutung für euch? Ich meine,
wie soll ich sagen, mir schien, als wäre ich mitten in eine Inszenierung
geraten. Musik, Campari und das grossartige Nachtessen, alles deutete
auf eine Feier hin?»

Joh und Manuela tauschten einen vielsagenden Blick aus. Joh griff nach
seiner Bierflasche und lehnte sich zurück. Louis war gespannt.
Über Manuela schien ein Energieschub zu kommen.

«Ja, wir haben gefeiert, der Tag gestern ist Teil unserer Rituale.
Gestern war Joh dran, nächste Woche ich. Jeden Monat einmal verwöhnen
wir uns gegenseitig während eines Tages.»

Manuela holte Luft.

«Ja, wir zelebrieren unser Zusammensein und überraschen uns mit
wertvollen, gemeinsamen Erinnerungen oder Visionen, was uns gerade einfällt, und manchmal steigen wir in den \emph{Hotpot} draussen. Und das Besondere daran, der Verwöhn-Tag ist nicht verschiebbar, findet statt, egal, was gerade um uns abgeht, und gelegentlich geht viel ab.»

Manuela hob die Augenbrauen. Bei Joh sah Louis ein Grinsen. Er wartete
gebannt auf eine Fortsetzung. Sie kam sofort.

«Seit zwei Jahren leben wir diese Entscheidung, wir mussten viel üben,
doch jetzt ist es erstaunlich, wie gut wir’s hinbekommen, nicht Joh?»

Der nickte und sah beflügelt aus.

«Gestern zum Beispiel war einiges los, David ist ausgerastet und du bist
unerwartet aufgetaucht, und trotzdem, wir halten an unserer Vereinbarung
fest. Das wissen auch unsere Arbeiter und Johs Jungs. Wir haben gelernt,
alles drumherum spontan einzubeziehen, doch mit dem Fokus auf uns zwei.
Das ist uns heilig. Und manchmal hauen wir einfach ab. Wie Teenager, die
die Zweisamkeit suchen. Es klappt immer besser.»

Manuela kicherte. Sie griff nach Johs Arm und drückte ihn. Die zwei schienen ein beachtliches Beziehungskonzept entwickelt zu haben. Louis war überrascht. Gleich ging es in ähnlicher
Manier weiter.

«Und wenn wir urlaubsreif sind, legen wir einen Zielort fest, steigen
mit einem Campari in die Hängematte und lassen «Campari Soda» laufen.
Und dann, dann kommt der wichtigste Teil,» Joh schaute verzückt vor sich
hin, «dann erzählen wir uns fiktive Abenteuergeschichten, die wir
zusammen erleben möchten, abwechselnd versteht sich, bis wir uns
treffen.»

Louis verstand nicht ganz. Aber es klang interessant.

«Und weisst du,» fügte Joh entspannt hinzu, «von Zeit zu Zeit geraten die
Dinge auch ausser Kontrolle und wir zoffen uns. Heilige sind wir weiss
Gott nicht.»

Johs Nachtrag holte Louis zurück auf den Boden. Er war beruhigt. Bald
wären sie ihm als Übermenschen entwischt.

«Nun aber Schluss mit unsern Geschichten, Ciro,» Manuela setzte sich
wieder, «magst du uns noch etwas singen? Einen Song, der dir am Herzen
liegt, etwas Spanisches vielleicht, als Halbspanier?»

Sie zwinkerte ihm zu.

Er brauchte etwas Zeit, um auf die Frage einzugehen. Ihm war, als hänge
er noch tief in ihren Erzählungen fest.

«Eh, ja, ehm, doch, mache ich sogar gerne.»

Louis nahm seine Gitarre in den Arm.

«Ja, ich überlege gerade, was euch gefallen könnte.»

«Egal, Ciro, \emph{dich} soll’s freuen, dann wird’s auch uns gefallen,»
meinte Joh und zeigte wieder sein charmantes Gesicht. Louis musste
unweigerlich lächeln.

«Also, wenn du einen spanischen Song wünschst, Manuela, dann spiel ich
einen, der mich durch meine Zeit des Abschieds von meiner Freundin
trägt, schon länger. Mir gefällt das Lied nicht nur musikalisch sehr
gut, auch der Text ist umwerfend.»

Manuela nickte mehrmals und schaute Louis mit grossen Augen an.
Eine Portion Offenheit würde ihm guttun. Er wollte es wagen, ihnen Einblick in einen traurigen Abschnitt seines
Lebens zu geben. Im Rückblick würde er sich über seinen Mut wundern.
Auch Joh nahm wieder Platz. Diesmal setzte er sich neben Louis. 

«Ja. Der Song ist von Ricardo Arjona und heisst 
\emph{«APNEA»,} Atemstillstand.»

Manuela und Johs  Aufmerksamkeit war geradezu greifbar. 

«Ich, ich habe mich in Arjonas Worten wiedergefunden. Mir gefällt, wie er die Folgen einer Trennung beschreibt, wie er singt - "

Manuela fuhr dazwischen.

« - unser Spanisch, Louis, ist nicht gerade grossartig,» und nach einem Seitenblick zu Joh sprach sie weiter, «kannst du den Inhalt vielleicht kurz zusammenfassen?»

«Ja, natürlich,» Louis sammelte sich, «warte, ehm, ich übersetze euch einige Sätze,» er schaute zur Decke hoch, « - in diesen Dienstagabend passt ein ganzes Jahrhundert. Der Spiegel wirft Pfeile des Vorwurfs. Heute beginnt das, was bereits vorbei ist. Die Hoffnung hat sich aus dem Fenster gestürzt, die Schlaflosigkeit ist hier her gezogen, um zu bleiben. Das Gestern hat alles auf morgen verschoben. Ich kann nicht atmen. Ich leide unter Atemstillstand, seit dem Tag, an dem du nicht mehr bei mir bist. Ich falle bis zum Grund des Meeres und kratze an der Blase, in der du nicht bist. Es ist unmöglich zu atmen, denn der Sauerstoff hat diesen Ort verlassen. - »

Joh räusperte sich und hielt seine Bierflasche mit beiden Händen
umschlossen.

«Nicht gerade leichte Kost.»

Er hob die Stirn, stand auf und setzte sich in eine bequeme Position neben Manuela.

Während seiner letzten Worte hatte Louis begonnen, die Gitarre zu stimmen. Er war soweit, 
setzte sich gerade auf und schloss für einen Moment die Augen. Der
Raum strahlte eine beachtliche Präsenz aus.
Dann setzte Louis seine Finger zum Spiel an. Raffiniert hatte er Arjonas
Version an seine Möglichkeiten als Gitarristen angepasst. Seinen
einführenden Text begleitete er mit einzelnen Tönen und sprach leise.
Den Song selbst steigerte er bis zur raumfüllenden Lautstärke. Verlor
sich hemmungslos in Klängen und Worten und fühlte, wie ein sich
steigerndes, wuchtiges Beben durch seinen Körper trieb.

\qrblock{Ricardo Arjona \emph{«APNEA»}}{media/QR-Codes/032_APNEA_Ricardo_Arjona_Spotify.png}

Louis’ Barriere brach zum Ende des Songs. Einer Dammsprengung gleich.
Jede Abwehr war wirkungslos. Sein Körper begann sich von oben nach unten
zu schütteln. Erst leicht, schliesslich mit einer Heftigkeit, die ihn
befremdete. Das Beben ging in ein Wimmern über. Es kam von tief unten.
Die Gitarre war bereits weggerutscht. Seine Hände vor dem Gesicht sass er mit gesenktem
Kopf da. Dass Joh und Manuela ihn in diesem erbärmlichen Zustand sahen,
war ihm nun vollkommen gleichgültig. Er schluchzte, und schniefte,
als wäre er hier alleine und fühlte sich wie ein überaltertes Kind. 
Die Zeit schien aufgelöst. Louis nahm leise Schritte wahr, die näher
kamen. Dann spürte er eine Hand, die ihn sanft am Kopf berührte. Er
blickte hoch und durch einen Schleier von Tränen sah er Joh vor sich.
Der lächelte ihn gütig an. Joh setzte sich langsam neben Louis, legte
ihm seine Hand auf die Schulter, drehte ihn zu sich um und schloss ihn 
in die Arme. Johs Zuwendung löste augenblicklich einen
weiteren Weinkrampf aus. Ein irritierendes Gefühl liess Louis kurz
erstarren, bis sich sein Widerstand ganz legte und er sich einfach
fallen liess. Es war zu schön, um falsch zu sein. Joh hielt ihn.

«Alles gut, Ciro, wir sind für dich da, wir haben Zeit.»

Louis hörte Johs Worte von weit weg.
Er wusste nicht, wie lange sie in dieser Umarmung dagesessen hatten.
Manuela schien jedenfalls geraume Zeit in der Küche herumhantiert zu
haben.
